\subsection{A.3. Divided powers on the weighted cohomology}

We use the graded divided-power convention: the algebra is strictly
alternating, and divided-power identities are imposed on even elements
of its augmentation ideal. We may denote higher divided powers of an
odd homogeneous element by zero, but do not impose the ungraded
addition identity on sums of odd elements. The differential identity
$d\gamma_m(z)=\gamma_{m-1}(z)dz$ is imposed for even $z$.
These distinctions matter in characteristic two. Scan 220 states
compatibility on $\Gamma^+M$; it does not impose an inconsistent
ungraded identity on odd exterior generators.

\paragraph{A-010. Proved: the canonical alternating PD algebra.}
For free modules $V,W$ over any ring $k$,
$\Gamma_k(V)\otimes_k\Lambda_k(W)$ has canonical graded divided
powers. For a positive even exterior form $a=\sum_I a_I e_I$,
$\gamma_m(a)$ is the sum over unordered collections of $m$ pairwise
disjoint nonempty even subsets $I$, with coefficient $\prod_I a_I$
and the sign of their wedge. In particular every positive decomposable
exterior monomial has higher divided powers zero. Degree-zero divided
powers are those of $\Gamma(V)$, and mixed even elements use the
addition identity.

\emph{Proof.}
Over the torsion-free universal ring on the finitely many coefficients,
the formula is $a^m/m!$: repeated exterior monomials vanish and the
$m!$ orderings of distinct even factors have the same sign.
For coefficient PD monomials the analogous integrality counts
partitions into unlabeled equal blocks. Thus the formulas are integral
polynomial formulas. The identities hold after rationalization and
hence in the torsion-free universal ring; arbitrary base change proves
them over $k$, including rings with $2$- and $3$-torsion. This proves
basis independence as well. Filtered colimits cover arbitrary free
ranks.

In $\Gamma\Phi\otimes\Lambda\Phi$, put $e=\sum_i dX_i$.
The ideal $(e)$ is differential and stable under even divided powers.
Indeed choose $e$ as a basis vector. Every exterior monomial in an
even element of this ideal contains $e$, so every product of two such
monomials is zero and all higher divided powers vanish. The quotient
is the simplicial PD-DGA $\Gamma\Phi\otimes\Lambda\Psi$.
Multiplication by $e$ identifies its degree-$q$ module with a submodule
in exterior degree $q+1$, but anticommutes with the differential and
is not an algebra embedding: products of two elements in its image
vanish.\hfill$\square$

\paragraph{A-011. Proved: descent to cohomology without factorial
cancellation.}
For every $X,k$, every even $q\geq0$, and every
$w\geq\max(1,q)$, divided powers induce
\[
 \gamma_m:H^q(C_w(X;k))\longrightarrow H^{mq}(C_{mw}(X;k)).
\]
They make the cohomology augmentation ideal a graded PD ideal. In
cohomological degree zero this ideal consists of positive-weight
coefficients; no higher divided powers are assigned to arbitrary
weight-zero elements of $k$.
The same assertion holds for literal locally finite weights and for
total-weight completion, interpreted by A-002 and A-005.

\emph{Proof.}
For $q>0$, the epimorphism $A_w^{q-1}\to Z_w^q$ is a Kan
fibration with contractible source, since $w-q+1>0$. Therefore
actual cocycles modulo actual boundaries are exactly $[X,Z_w^q]$.
A boundary factors through a contractible source; conversely a
nullhomotopy lifts starting at zero and produces an actual primitive.
For even $q$, $\gamma_m$ is a simplicial set map
$Z_w^q\to Z_{mw}^{mq}$, because its differential is zero on cycles.
A simplicial map preserves homotopies, hence induces the stated
operation. Every PD identity descends because it already holds on
forms. In degree zero there are no boundaries. The same argument
uses the weight-preserving contractions on the direct sum, or product,
of all weights at least $q$. The operations are respectively locally
finite or convergent. No integer is canceled in this proof.
\hfill$\square$

\paragraph{A-012. Proved: an explicit integral formula for every
induced operation.}
The following finite algorithm completely determines the operation
on ordinary cocycles, including all torsion cases. Fix even $q>0$,
$w\geq q$, $m\geq1$, and $N=mq$. For every $N$-simplex $\sigma$
and ordinary normalized $q$-cocycle $z$, restrict $z$ to $\Delta[N]$.
For $1\leq c_1<\cdots<c_q\leq N$, let
$\alpha_c:[N]\twoheadrightarrow[q]$ have blocks
\[
 [0,c_1-1],\ [c_1,c_2-1],\ \ldots,\ [c_q,N].
\]
Write $\varepsilon_c=\alpha_c^*\varepsilon_q$, where
$\varepsilon_q$ is the orientation cocycle on $\Delta[q]$.
Solve $z|_\sigma=\sum_c u_c\varepsilon_c$ by evaluation on the
faces $(0,j_1,\ldots,j_q)$, ordered compatibly with componentwise
order. The matrix entry is one exactly when
$c_r\leq j_r<c_{r+1}$ for $r<q$ and $c_q\leq j_q$.
Its diagonal entries are one and all nonzero entries have $c\leq j$;
thus the solution uses only addition and subtraction in $k$.

Put $U_{c,r}=\sum_{\alpha_c(v)=r}X_v$ and form
\[
 L_{w,q}(z|_\sigma)=\sum_c u_c U_{c,0}^{[w-q]}
                           dU_{c,1}\wedge\cdots\wedge dU_{c,q}.
 \tag{A-10}
\]
Expand $\gamma_m$ of (A-10), replace $dX_0$ by
$-\sum_{v=1}^N dX_v$, and take the coefficient of
$dX_1\wedge\cdots\wedge dX_N$. Replace every PD monomial
$X_0^{[a_0]}\cdots X_N^{[a_N]}$ in it by one. The resulting value
is $P_{m,w,q}(z)(\sigma)$. It is a normalized $mq$-cocycle and
\[
 \gamma_m([z]T^{[w]})=[P_{m,w,q}(z)]T^{[mw]}.
 \tag{A-11}
\]
Degree zero is treated below, odd higher graded powers are zero,
and sums of even classes in different weights use the addition
formula. For literal weights, an ordinary $I_q$-valued cocycle need
not be a globally finite sum of weight classes. On each simplex,
apply the displayed lift to the finitely many weight components
present on its finitely many faces, sum those lifts, and then apply
$\gamma_m$. The same coefficient extraction gives an
$I_{mq}$-valued cocycle. For total-weight completion, perform this
construction one output weight at a time; only finitely many input
weights contribute. Thus the rule also specifies every induced
operation in these two conventions.

\emph{Proof.}
The normalized simplex form
$\omega=X_0^{[w-q]}dX_1\wedge\cdots\wedge dX_q$ is a cycle
of $NZ_w^q$ in simplicial degree $q$: every boundary has dimension
$q-1$. Its integral is $T^{[w]}$. Sending the generator of $k[q]$
to $\omega$ is a chain map; Dold--Kan extends it to a simplicial
linear map $K(k,q)\to Z_w^q$. Its degeneracy decomposition is
exactly (A-10). To see the stated basis directly, a cocycle on
$\Delta[N]$ is determined by its values on faces containing vertex
zero (the cocycle equation determines every other face), and the
displayed unitriangular matrix has determinant one. Both source
and target have sole homotopy group $k$ in degree $q$, and the
integral makes the map the identity there. Hence it induces the
inverse weightwise comparison on maps out of every simplicial set.

Integration on $\Delta[N]$ sends
$X_0^{[a_0]}\cdots X_N^{[a_N]}dX_1\cdots dX_N$ to
$T^{[N+\sum a_v]}$. Compute over $\mathbb Q$ by the Dirichlet
simplex integral: the factorials cancel. Both sides are integral in
the universal torsion-free PD algebra, so the identity and Stokes'
identity hold integrally and survive arbitrary base change. This
proves that coefficient extraction is exactly the canonical cocycle
comparison. A-011 proves independence of the cocycle representative. For the
literal extension, Dold--Kan gives
$K(I_q,q)=\bigoplus_{w\geq q}K(k,q)$, because in each simplicial
degree its defining sum is finite. The direct sum of the specified
maps $K(k,q)\to Z_w^q$ is a weak equivalence onto
$\bigoplus_{w\geq q}Z_w^q$. It therefore gives exactly the
coefficient comparison with $H^q(X;I_q)$ of A-002. The product of
these maps gives the corresponding completed comparison. These
arguments do not decompose an arbitrary literal cohomology class
into globally finitely many weight classes.
\hfill$\square$

\paragraph{A-013. Proved under the stated hypotheses: useful scalar
formulas and their limits.}
For $m\geq1$ and $w>0$, set
$c(m,w)=(mw)!/(m!(w!)^m)$, an integer. In degree zero, on the
positive-weight PD ideal,
\[
 \gamma_m(xT^{[w]})=c(m,w)x^mT^{[mw]}.
\]
The same formula holds for an even positive-degree class whenever
multiplication by $m!$ is injective on its target cohomology group.
Indeed the PD identity $m!\gamma_m(a)=a^m$ and A-003 give the
formula after multiplying by $m!$, and the stated injectivity allows
cancellation. Without it, A-012 is the unconditional answer; we do
not infer equality by an invalid torsion cancellation. For example,
the formula applies to the integral even classes of $\mathbb{CP}^2$.
An unqualified formula for odd classes is false: on
$\mathbb{RP}^2$ over $\mathbb F_2$, the degree-one class $x$ has
$x^2\ne0$. In weight one, the unqualified scalar formula would
give $\gamma_2([x]_1)=[x^2]_2\ne0$, whereas its higher graded
divided power is zero by the odd-degree convention.
The finite bar-cochain verification is included with B-09.

